% Section 6: Case Study
% Source: context/Case_Study_OpenOrderbook.md (all sections §CS1-§CS9)
% Source: context/Paper_Layout_Plan.md §6
%
% EM-DASH PROHIBITION: Do not use --- or \textemdash. Use colons or semicolons.

\section{Industrial Case Study}\label{sec:case_study}

This section addresses \RQ{5} using a mixed-method design: qualitative
triangulation across five evidence sources is augmented by the ecological
validity arm as an embedded controlled experiment (the no-Vault counterfactual).

\subsection{Case Context}\label{sec:case_context}

% ANON: All identifying details below are anonymised for double-anonymous review.
% Camera-ready version will name the company, product, and personnel with consent.

The case is a
\anon{production blockchain trading platform}{production OpenOrderbook/Upstream Exchange platform}
operated by
\anon{a Swiss-based fintech company}{Abbey Technologies GmbH}
serving approximately 13{,}000 registered users.
% ANON: company name, product name, exchange name
The platform provides fixed-return options trading in tokenised financial
instruments through an AI-native interface: an LLM orchestrator translates
trader instructions into on-chain operations via MCP tool chains. The
blockchain infrastructure is a private Proof-of-Authority (PoA) Ethereum-compatible
network with four active nodes.

Authentication uses RS256~\cite{rfc7518} JWTs~\cite{rfc7519} issued by Azure AD B2C at approximately
1{,}080 characters, placing the production tokens in the \emph{very long} tier
of our experimental design (Section~\ref{sec:objects}). A single character
corruption invalidates the RSA signature and produces an HTTP~403
rejection at the downstream API~\cite{rfc6750}. The LLM configuration supports multiple
providers (the controlled study includes Claude, GPT, and Gemini families);
a local privacy-preserving mode using Gemma~4~E4B (4B-parameter open-weight,
Ollama) is also supported for data-sovereignty deployments.

\textbf{Concrete instantiation.}
\texttt{[VP:~LLM]} = user-configurable model; \texttt{[VP:~Protected Resource]} = PoA
Ethereum RPC + REST trading API; \texttt{[VP:~Credential-Bearing Tool]} = RS256
signing service issuing Azure AD B2C JWTs.

\subsection{Evidence Sources}\label{sec:evidence}

Five independent evidence sources are triangulated. (1)~Production source code
and git commit history provide an objective timeline of mitigations.
(2)~Contemporaneous developer communications document the misdiagnosis phase
and the qualitative observation that prompting ``gave better results.''
(3)~A structured interview with
\anon{the company's CEO and infrastructure engineer}{Brian Collins and Pete Hall}
% ANON: personnel names
provides architectural context and co-author fact verification.
(4)~The ecological validity arm ($N\!=\!1{,}600$) serves a dual role: it
quantifies the production-specific \DFR{} for the main study, and it
constitutes the controlled counterfactual strand of this case study---the
\emph{no-Vault} condition instantiated on the actual production system via
an environment variable flag (\texttt{RESEARCH\_RAW\_MODE=true}), providing
a quantified baseline against which the post-Vault observation is interpreted. (5)~Commit-level pre/post evidence confirms zero
authentication failures after Vault deployment. No source contradicts another.

\subsection{Discovery Trajectory}\label{sec:trajectory}

JWT corruption was first observed on 2026-04-03 during UAT testing, eight days
after the MCP server was first built. The observable symptom was intermittent
\texttt{HTTP 403 Unauthorized} responses from the downstream REST API on
authenticated trading operations. The intermittent failure rate was
approximately 20\%, consistent with the very-long tier \DFR{} in the
controlled study.

\textbf{Misdiagnosis.}
The failure was attributed to infrastructure: signatures failed
\texttt{RSA.VerifyData()} against all six active JWKS keys. One key was 47 days
old, leading to the hypothesis that Azure AD was issuing tokens with a stale
private key; a Microsoft Support case was raised. The misdiagnosis was plausible:
a JWT corrupted by single-character substitution is cryptographically
indistinguishable from a token signed with an invalid key.

\textbf{Root cause identification.}
\anon{The lead developer}{Andrew Le Gear},
% ANON: personnel name
who had personally observed the corruption during testing, independently
formed the LLM relay hypothesis from knowledge of how LLMs handle opaque,
high-entropy strings. A diagnostic mechanism was added to the tool chain: the
\texttt{acquire\_bearer\_token} tool was modified to return a
\texttt{tokenFingerprint} (SHA256[:16] of the JWT), and a verification tool
was added to compute the fingerprint of the received bearer token. Comparing
fingerprints before and after LLM relay confirmed that the token was being
altered in transit.

\textbf{Prompting-first mitigation.}
The initial response followed a prompting-first trajectory: a 318-character
\texttt{CRITICAL} verbatim fidelity instruction was added to the
\texttt{bearerToken} parameter in the tool registry. The instruction explicitly
names the RSA signature consequence and the single-character precision
requirement.
The partial improvement confirmed the relay hypothesis but did not solve the
problem. This is precisely the intervention evaluated in the ecological
validity arm (Section~\ref{sec:eco_results}): the same 318-character
instruction, the same production JWT schema, producing 19.9\% residual \DFR{}.

\textbf{Architectural solution.}
On 2026-04-13 (commit \texttt{5df1b265a}),
\anon{the lead developer}{Le Gear}
% ANON: personnel name
deployed the Vault implementation, with the commit comment: \emph{``Bearer token
vault: bypass the LLM to prevent token corruption. The LLM hallucinates data
into JWT tokens (expanding claims, appending fake responses). We intercept the
token from acquire\_bearer\_token results and inject it directly.''}
This comment identifies two failure modes observed in production:
(1)~semantic expansion, where the LLM hallucinated additional JSON claims
into the JWT payload; and (2)~content fabrication, where the LLM appended
fabricated tool-response content after the JWT. Both are consistent with the
failure taxonomy in Table~\ref{tab:taxonomy}.

The commit message uses no architectural pattern vocabulary; the solution was
derived from first principles under production pressure. Pattern identification
and naming were applied retrospectively, confirming the RA as a naturally
emerging solution that this paper names, formalises, and validates.

\subsection{Reactive Mitigations as Defence-in-Depth}\label{sec:reactive}

Before the Vault was implemented, and retained alongside it as defence-in-depth,
two reactive mitigation layers were added to the MCP handler. The first,
\texttt{SanitizeBearerToken()}, handles structural failures: it extracts the
\texttt{accessToken} field when the LLM passes the full JSON response object
(case~1), strips non-Base64url characters appended after the signature
(case~2), and removes \texttt{Bearer} prefix confusion (case~3). The second,
\texttt{JwtValidator.ValidateAndRepairAsync()}, implements single-character
brute-force repair of RS256 signatures: for each character position in the
signature segment, it substitutes all 64 Base64url characters and tests
\texttt{RSA.VerifyData()} against the JWKS keys. This repairs Levenshtein-1
corruption in the signature segment but cannot repair multi-character corruption
or payload-level semantic substitution.

Table~\ref{tab:coverage} maps these layers against the failure taxonomy. The
reactive mitigations cover format errors and low-Levenshtein signature
corruption but cannot detect semantic substitution (\texttt{sem\_sub}), where
the corrupted token is well-formed and plausible-length. The Vault is the only
mechanism with complete coverage: it bypasses the failure mode entirely by
removing the credential from the LLM's inference path.

\begin{table}[t]
\centering
\footnotesize
\vspace{-4pt}\caption{Defence-in-depth coverage matrix. Reactive mitigations cover a subset
of failure modes; the Vault covers all by bypassing LLM relay.}
\label{tab:coverage}
\begin{tabular}{@{}lccc@{}}
\toprule
\textbf{Failure mode} & \textbf{Sanitize} & \textbf{Validator} & \textbf{Vault} \\
\midrule
JSON object relay        & \checkmark & Partial & \checkmark \\
Trailing fabrication     & \checkmark & Partial & \checkmark \\
Bearer prefix confusion  & \checkmark & Partial & \checkmark \\
Char-level substitution  & $\times$   & \checkmark & \checkmark \\
Multi-char corruption    & $\times$   & $\times$   & \checkmark \\
Abstention / refusal     & $\times$   & $\times$   & \checkmark \\
\bottomrule
\end{tabular}
\end{table}

\subsection{Ecological Validity}\label{sec:eco_results}

The eco arm ($N\!=\!1{,}600$; 16 models $\times$ 100 iterations) is the
controlled counterfactual for this case study. It used the exact production JWT
schema (1{,}080-character RS256 Azure AD B2C token) and the verbatim
318-character \texttt{BearerTokenParam} instruction from the production
codebase, executed through the production software with Vault protection
deliberately disabled. An environment variable flag (\texttt{RESEARCH\_RAW\_MODE=true}, commit
\texttt{2396d5273}, 2026-07-07) bypasses both \texttt{SanitizeBearerToken()}
and \texttt{JwtValidator}, forwarding the raw LLM output verbatim to the
downstream API. The production binary is identical in both modes; only the
runtime environment variable differs, ensuring the treatment and control
conditions share the same code path and differ solely in the presence or
absence of Vault interception. This design was necessary because the
production mitigations are effective enough to mask the underlying failure:
without deliberately disabling them, the raw \DFR{} is not observable.

The eco arm \DFR{} was 19.9\%, a $+1.9$ percentage point difference from the
main study ($p\!=\!0.16$, Wilcoxon; not
significant). The Spearman rank correlation between eco arm and main experiment
model rankings was $\rho\!=\!0.740$ ($p\!=\!0.001$), confirming that model
rankings transfer to production conditions.

\subsection{Post-Vault Deployment}\label{sec:post_vault}

Zero JWT authentication failures have been recorded over the 88-day observation
window from Vault deployment (2026-04-13) through 2026-07-10. Azure Monitor
metrics record 4{,}774 MCP tool invocations during this window; at the eco
arm's 19.9\% \DFR{}, approximately 950 would have failed without Vault
protection. The eco arm (Section~\ref{sec:eco_results}) provides the quantified
counterfactual on the same production system.
The post-Vault zero-failure observation is interpreted against that baseline,
not in isolation. This is a structural result, not a statistical one: the Vault
eliminates the failure mode by removing the credential from the LLM's inference
path. The expected post-Vault \DFR{} is not ``lower than 18.0\%'' but exactly
zero by construction for the byte-relay property: the byte sequence injected
into each outbound request is identical to the credential stored in the Vault.
This guarantee does not extend to credential theft, token expiry, revocation,
or replay attacks, which require independent controls.
